\documentclass[10pt,a4paper]{amsart}
\usepackage[margin=19mm]{geometry}
\usepackage{amsmath,amssymb,mathtools}
\usepackage[scr=rsfs,cal=euler]{mathalfa}
\usepackage{xfrac}
\usepackage[unicode,hidelinks,bookmarks=false]{hyperref}
\newcommand{\ie}{i.e.\ }
\newcommand{\cf}{cf.\ }
\newcommand{\resp}{respectively\ }
\newcommand{\Subsection}{\subsection}
\providecommand{\nobreakdash}{\nobreak\hbox{-}}
\setlength{\emergencystretch}{2em}
\multlinegap0pt
\usepackage{amssymb}
\usepackage{xcolor}
\newtheorem{lemma}{Lemma}
\newtheorem{thm}{Theorem}
\title{Rings in which one-sided strongly $\pi$-regular elements are strongly $\pi$-regular}
\author{Dimple Rani Goyal, Dinesh Khurana}
\date{Source: arXiv:2511.20070v2 (CC BY 4.0)}
\begin{document}
\maketitle

\noindent\emph{Source and licence.} Rings in which one-sided strongly $\pi$-regular elements are strongly $\pi$-regular. Version arXiv:2511.20070v2 (CC BY 4.0), distributed by arXiv under the Creative Commons Attribution 4.0 International licence, http://creativecommons.org/licenses/by/4.0/, as recorded in the arXiv licence metadata for this exact version and shown on the abstract page https://arxiv.org/abs/2511.20070v2. Full licence notice: \texttt{licenses/arxiv-2511.20070-CC-BY-4.0.txt}.\par\smallskip

\noindent\textit{Editorial scope.} Counted unit: the article's thm D.mainthm (source0.tex lines 276--278). The article's complete original proof of it is delivered with it (source0.tex lines 279--289, one \texttt{proof} environment of 11 lines). No mathematics is written, repaired or re-derived: the statement and the proof are the article's own lines, in the article's own notation, and no retained line is substituted or re-flowed.  Retained uncounted as the source-local context the proof needs: lines 272--274.  Nothing was omitted from any retained span, so the package carries no omission record.  No retained line contains a citation, so the wrapper carries no bibliography and no bibliography entry.  No retained line contains a figure environment, an image-inclusion command or drawing code, so no image is delivered and the package declares no asset dependency.  Editorial formatting only, in addition: an AMS \texttt{amsart} preamble that restates the article's own declarations and macros used by the retained lines, namely the environments \texttt{lemma}, \texttt{thm} on separate counters, together with the standard packages those lines need.  The retained environments print in this document as Lemma 1, Theorem 1 in order of appearance, whereas the article prints the same items with the numbering of its own full document.  The complete native source of the article as distributed in the arXiv e-print for this exact version is bundled unchanged as \texttt{sources/189761/source0.tex}.
	\begin{lemma}\label{D.mainlemma}
		If $m_1$ and $m_2$ are monomials in $\mathcal{B}$, then either $m_1m_2$ is already reduced or $m_1m_2 \prec m_2$ after reduction. 
	\end{lemma} 

	\begin{thm}\label{D.mainthm}
		If $f \in R$ has a monomial in its support that is in $Rx$, then $r_R(f) = 0$. In particular,  $r_R(g) = 0$ for any $g \in R$ that has $1$ in its support.
	\end{thm}

	\begin{proof}
		Suppose $h$ is any nonzero element in $R$. We will show that $fh \neq 0$.\\
		Let $M = x^{i_{n}}a\cdots x^{i_{2}}ax^{i_1}a^{i_0}$ be the maximal monomial in the support of $h$, and $m$ be a monomial of the largest length in the support of $f$ that is in $Rx$. 
		
		If $mM$ appears only once in the support of $fh$, then we have nothing to prove. If there exist monomials $m_1$ in the support of $f$ and $m_2$ in the support of $h$ such that $mM = m_1 m_2$, then $m_1m_2$ is irreducible because otherwise, by Lemma \ref{D.mainlemma}, $mM = m_1m_2 \prec m_2 \prec M$, a contradiction. Due to the maximality of $M$, the monomial $m_2$ is a right subword of $M$. Also $m_2 \prec M$ because otherwise $mM = m_1M$ implies that $m = m_1$. So the length of $m_1$ is greater than that of $m$, and as $m$ is a monomial of the largest length in $Rx$,  $m_1 \in  Ra$. This implies $m_2 = M(k)$ for some $k \in \{1, \ldots, n\}$. Let $l \in \{1,\ldots, n\}$ be the smallest integer such that $M(l)$ is in the support of $h$. We will show that $mM(l)$ appears in the support of $fh$ only once.
		
		Suppose $mM(l) = m_3m_4$ for some monomial $m_3$ in the support of $f$ and $m_4$ in the support of $h$ such that $(m_3, m_4) \neq (m, M(l))$. First note that $m_3m_4$ is irreducible because otherwise, as $m \in Rx$, we have $M \prec mM(l) = m_3m_4  \prec  m_4$ by Lemma \ref{D.mainlemma}, a contradiction. If $m_3 = mm_0$ for some monomial $m_0$, then $m_0  \in Ra$ as $m$ is the largest monomial in the support of $f$ that is in $Rx$. But then $m_4 = M(t)$ for some $t < l$, violating the minimality of $M(l)$. Lastly, if $m_4 = m_0M(l)$ for some monomial $m_0$, then as $mM(l) \in RxM(l)$, so $m_0 \in Rx$. But this implies  $M \prec m_4$,  a contradiction.  
		So $mM(l)$ is a nonzero monomial in the support of $fh$ and so $fh \neq 0$. So $r_{ R}(f)=0$.
		
		For the last part, if $g \in R$ has $1$ in its support, then $xg$ has $x$ in its support. For any element $g_1 \in R$ if $gg_1 = 0$ then $xgg_1 = 0$ and by the first part, $g_1 = 0$. Thus $r_R(g) = 0$. 
	\end{proof} 

\end{document}
